% CPEN 442 -- Assignment 2 answer template (optional).
% Fill in your name(s) and answers below, compile, and submit the PDF as a2.pdf.
% You do not have to use LaTeX: any readable PDF is fine.
\documentclass[12pt]{article}

\usepackage[margin=1in]{geometry}
\usepackage[dvipsnames,table]{xcolor}
\usepackage[most]{tcolorbox}
\usepackage{hyperref}
\usepackage{listings}
\usepackage{graphicx}
\usepackage{amsmath}
\usepackage{tabularx}
\usepackage[T1]{fontenc}

\definecolor{alicepink}{rgb}{0.98,0.51,1.}
\definecolor{aliceblue}{rgb}{0.7, 1, 1}
\definecolor{background}{rgb}{1, 0.98, 0.95}

\lstset{frame=tb, language=Python, showstringspaces=false, columns=flexible,
  basicstyle={\small\ttfamily}, breaklines=true, backgroundcolor=\color{background}}

\setlength{\parindent}{0in}
\setlength{\parskip}{\baselineskip}

\newcounter{question}
\newcommand{\question}[1]{\refstepcounter{question}%
  \textbf{\colorbox{aliceblue}{[#1 pts]} Question \arabic{question}}}

% Box to write each answer in. It grows with its content.
\newtcolorbox{answer}{colback=white, colframe=black, boxrule=0.4pt,
  sharp corners, breakable}

\begin{document}

\begin{center}
    THE UNIVERSITY OF BRITISH COLUMBIA\\
    Electrical and Computer Engineering Department\\[\baselineskip]
    {\bf CPEN 442\hfill Introduction to Cybersecurity\hfill Winter Term 1, 2026}
    \vspace{1cm}

    {\sc \Large Assignment 2: Symmetric Crypto}
\end{center}

% One or two students. Delete the second row if you worked alone.
\begin{tabularx}{\textwidth}{|X|l|}
    \hline
    \rowcolor{gray!15} \textbf{Name} & \textbf{Student number} \\
    \hline
    First Last & 12345678 \\
    \hline
    First Last & 12345678 \\
    \hline
\end{tabularx}

Questions 4--7 are answered in \texttt{attacker.py}; this PDF covers Questions 1--3.

\section{Detecting Modes of Operation}

\question{15} Which of \texttt{lyrics-ciphertext1.txt}, \texttt{lyrics-ciphertext2.txt}, and \texttt{lyrics-ciphertext3.txt} is AES-ECB, and how did you tell?

\begin{answer}
    \textbf{AES-ECB is:} \texttt{lyrics-ciphertext?.txt}

    \textbf{Explanation:}
    % Write your explanation here.
\end{answer}

\question{15} Which one is AES-CTR, and how did you tell?

\begin{answer}
    \textbf{AES-CTR is:} \texttt{lyrics-ciphertext?.txt}

    \textbf{Explanation:}
    % Write your explanation here.
\end{answer}

\question{30} Give the mode of operation of each image in \texttt{ciphertext-sprites}, and explain briefly how you found it.

\begin{answer}
    % You can include code with \begin{lstlisting} ... \end{lstlisting}
    % or figures with \includegraphics in any of the explanations.
    \textbf{AES-ECB:} ?, ?, ?, ?

    \textbf{Explanation:}
    % Write how you identified the AES-ECB images here.

    \bigskip
    \textbf{AES-CBC:} ?, ?, ?, ?

    \textbf{Explanation:}
    % Write how you identified the AES-CBC images here.

    \bigskip
    \textbf{AES-CTR:} ?, ?, ?, ?

    \textbf{Explanation:}
    % Write how you identified the AES-CTR images here.
\end{answer}

% Include this section only if you used generative AI.
% It must be between half a page and one page long, and cannot be written by AI.
\section*{Generative AI Disclosure}

% Describe how and where you used generative AI.

\end{document}
